\begin{frame}[t]{\ev{\tagP\,\tagO}Promotion is small enough to model-check}
\lead{A TLA+ design model of one cell, from epoch replacement to publication.}
\vspace{-.15cm}
{\small
\begin{ticks}
\item evidence refers only to the frozen candidate
\item approval requires evidence for that same candidate
\item publication requires frozen candidate $=$ evidence $=$ approval under the current Airflow-bound epoch
\item an external commit belongs to the current epoch
\item a search-origin authority request cannot become an authority source
\item an \texttt{in\_doubt} effect has no blind retry and leaves that state only by observation
\end{ticks}\par}
\vspace{.05cm}
{\small\color{Rust}\bfseries Design model: no TLC model-checker run and no trace check recorded yet. The epoch fence is proposed.\par}
\vspace{.1cm}
\colorbox{Pale}{\parbox{\dimexpr\textwidth-2\fboxsep\relax}{\footnotesize Dean \& Barroso note that tail-tolerant techniques suit operations without critical mutations, while consistent updates use quorum algorithms such as Paxos.}}
\claim{Model-check the gate and count the children.}
\sources{\href{https://github.com/zozo123/ariflow-swfactory/tree/main/formal/authority}{TLA+ authority model of our loop}. Dean \& Barroso, The tail at scale, CACM 2013 (``Mutations'').}
\note{[0:45] The promotion gate must never race. It is small enough for a tool to model-check by trying every ordering. Evidence names only the frozen candidate. Approval needs that evidence. Publication needs all three in the current epoch. Six rules must always hold, and no model-checker run exists yet. Dean and Barroso report that consistent updates use quorum protocols such as Paxos. A child's output cannot be model-checked. It must instead be counted.}
\end{frame}
